\documentclass[11pt]{article}

\usepackage[margin=1.7cm]{geometry}
\usepackage{amsmath}
\usepackage{amssymb}
\usepackage{enumitem}
\usepackage{titlesec}
\usepackage{xcolor}
\usepackage{fancyhdr}
\usepackage{array}
\usepackage{makecell}

\titleformat{\section}{\large\bfseries}{\thesection}{0.6em}{}
\titlespacing{\section}{0pt}{1.1em}{0.5em}

\pagestyle{fancy}
\fancyhf{}
\fancyhead[L]{Relative Risk from Two-Way Tables}
\fancyhead[R]{Name: \makebox[4cm]{\hrulefill}}
\fancyfoot[C]{\thepage}
\renewcommand{\headrulewidth}{0.4pt}

\newcommand{\ans}{\hrulefill}
\newcommand{\anshere}{\par\vspace{0.3em}\noindent\ans\par\vspace{0.3em}}
\setlist[enumerate,1]{leftmargin=1.4em, itemsep=0.5em, topsep=0.2em}
\setlist[enumerate,2]{leftmargin=1.3em, itemsep=0.3em, topsep=0.2em}

\newcolumntype{Y}{>{\centering\arraybackslash}p{2.1cm}}
\newcolumntype{L}{>{\raggedright\arraybackslash}p{3.1cm}}

\begin{document}

\begin{center}
  {\Large\bfseries Relative Risk from Two-Way Tables}
\end{center}

\vspace{0.3em}
\noindent A \textbf{two-way table} sorts data by two categories at once --- one category down the rows, one across the columns. Each row and column has a \textbf{total}, and the bottom-right cell is the \textbf{grand total} (everyone in the data set).

\vspace{0.5em}
\section{Two-Way Tables: The Basics}
\noindent Complete each table, then answer the questions.

\subsection*{1. Favourite subject}
\noindent 220 students were asked their favourite subject out of Maths, Art and PE.

\begin{center}
\begin{tabular}{|L|Y|Y|Y|Y|}
\hline
 & Maths & Art & PE & \textbf{Total} \\\hline
\textbf{Year 9}  & 45 & 30 & 25 & \\\hline
\textbf{Year 10} & 35 & 40 & 45 & \\\hline
\textbf{Total}   &    &    &    & \\\hline
\end{tabular}
\end{center}

\begin{enumerate}
  \item How many Year 10 students chose Art? \ans
  \item What fraction of all 220 students are in Year 9? \ans
  \item What percentage of Year 9 students chose PE? \ans
  \item Of all the students who chose Maths, what fraction are in Year 10? \ans
\end{enumerate}

\subsection*{2. Pet ownership}
\noindent Households in an urban and a rural area were asked what pets they own.

\begin{center}
\begin{tabular}{|L|Y|Y|Y|Y|}
\hline
 & Dog & Cat & No pet & \textbf{Total} \\\hline
\textbf{Urban} & 20 & 35 & 45 & \\\hline
\textbf{Rural} & 50 & 15 & 10 & \\\hline
\textbf{Total} &    &    &    & \\\hline
\end{tabular}
\end{center}

\begin{enumerate}
  \item How many rural households own a cat? \ans
  \item What fraction of urban households have no pet? \ans
  \item What percentage of all households own a dog? \ans
  \item Of all the households that own a cat, what fraction are rural? \ans
\end{enumerate}

\newpage
\section{Relative Risk}
\noindent In a two-way table about risk, the rows are usually two \textbf{groups} (e.g.\ exposed/not exposed, treated/untreated) and the columns are whether an \textbf{outcome} happened or not.
\[
  \text{risk (for a group)} = \dfrac{\text{number in that group \emph{with} the outcome}}{\text{total number in that group}}, \qquad
  \text{RR} = \dfrac{\text{risk in Group A}}{\text{risk in Group B}}.
\]
\noindent Note: risk uses the \textbf{row} total for that group, \emph{not} the grand total. \quad
If $RR=1$ the risk is the same in both groups; if $RR>1$ Group A has the higher risk; if $RR<1$ Group A has the lower risk.

\vspace{0.3em}
\noindent For each table below: complete the table, then find the risk for each group, the relative risk, and interpret it.

\subsection*{1. Vaccination and flu}
\noindent 1000 people took part in a study. 500 received a flu vaccine and 500 did not. The table shows how many of each group caught the flu that winter.

\begin{center}
\begin{tabular}{|L|Y|Y|Y|}
\hline
 & Flu & No flu & \textbf{Total} \\\hline
\textbf{Vaccinated}   & 40  &     & 500 \\\hline
\textbf{Unvaccinated} & 100 &     & 500 \\\hline
\textbf{Total}        &     &     &     \\\hline
\end{tabular}
\end{center}

\begin{enumerate}
  \item Risk of flu for the vaccinated group. \ans
  \item Risk of flu for the unvaccinated group. \ans
  \item Relative risk of flu for the vaccinated group compared to the unvaccinated group. \ans
  \item Interpret your answer to part (c) in a sentence.\anshere
\end{enumerate}

\subsection*{2. Energy drinks and sleep problems}
\noindent A study of 500 teenagers looked at whether regularly drinking energy drinks is linked to sleep problems.

\begin{center}
\begin{tabular}{|L|Y|Y|Y|}
\hline
 & \makecell{Sleep\\problems} & \makecell{No sleep\\problems} & \textbf{Total} \\\hline
\textbf{Regular energy drink user} & 72 & & 200 \\\hline
\textbf{Not a regular user}        & 36 & & 300 \\\hline
\textbf{Total}                     &    & &     \\\hline
\end{tabular}
\end{center}

\begin{enumerate}
  \item Risk of sleep problems for regular energy drink users. \ans
  \item Risk of sleep problems for people who are not regular users. \ans
  \item Relative risk of sleep problems for regular users compared to non-users. \ans
  \item Interpret your answer to part (c) in a sentence.\anshere
\end{enumerate}

\newpage
\subsection*{3. Seatbelts and injury}
\noindent Police records of 1000 car accidents were checked to see whether wearing a seatbelt made a difference to whether a person was injured.

\begin{center}
\begin{tabular}{|L|Y|Y|Y|}
\hline
 & Injured & Not injured & \textbf{Total} \\\hline
\textbf{Wearing a seatbelt}     & 30 & & 600 \\\hline
\textbf{Not wearing a seatbelt} & 80 & & 400 \\\hline
\textbf{Total}                  &    & &     \\\hline
\end{tabular}
\end{center}

\begin{enumerate}
  \item Risk of injury for people wearing a seatbelt. \ans
  \item Risk of injury for people not wearing a seatbelt. \ans
  \item Relative risk of injury for seatbelt wearers compared to non-wearers. \ans
  \item Interpret your answer to part (c) in a sentence.\anshere
\end{enumerate}

\subsection*{4. Exercise and heart disease}
\noindent 740 adults over 50 were tracked over several years to see whether regular exercise was linked to heart disease.

\begin{center}
\begin{tabular}{|L|Y|Y|Y|}
\hline
 & \makecell{Heart\\disease} & \makecell{No heart\\disease} & \textbf{Total} \\\hline
\textbf{Regular exercise}    & 25 & & 500 \\\hline
\textbf{No regular exercise} & 60 & & 240 \\\hline
\textbf{Total}               &    & &     \\\hline
\end{tabular}
\end{center}

\begin{enumerate}
  \item Risk of heart disease for people who exercise regularly. \ans
  \item Risk of heart disease for people who do not exercise regularly. \ans
  \item Relative risk of heart disease for the regular-exercise group compared to the no-exercise group. \ans
  \item Interpret your answer to part (c) in a sentence.\anshere
\end{enumerate}

\newpage
\subsection*{5. Sunscreen and sunburn}
\noindent 400 beachgoers were surveyed about sunscreen use and whether they got sunburnt that day.

\begin{center}
\begin{tabular}{|L|Y|Y|Y|}
\hline
 & Sunburn & No sunburn & \textbf{Total} \\\hline
\textbf{Used sunscreen}    & 15 & & 200 \\\hline
\textbf{Did not use sunscreen} & 30 & & 200 \\\hline
\textbf{Total}             &    & &     \\\hline
\end{tabular}
\end{center}

\begin{enumerate}
  \item Risk of sunburn for people who used sunscreen. \ans
  \item Risk of sunburn for people who did not use sunscreen. \ans
  \item Relative risk of sunburn for sunscreen users compared to non-users. \ans
  \item Interpret your answer to part (c) in a sentence.\anshere
\end{enumerate}

\vspace{0.5em}
\noindent \textbf{Reflection:} Which of the five tables above show a \emph{protective} factor (one that lowers risk), and which show a factor that \emph{increases} risk? \anshere

\end{document}
